\subsection{Ảnh hưởng của chế độ fine-tuning}

Ở cả hai kiến trúc, accuracy tăng khi cho phép nhiều tầng pretrained được cập nhật hơn. Với ResNet18, head-only đạt 60,89\%, còn full fine-tuning đạt 78,93\%. Với Swin-T, hai mức tương ứng là 71,75\% và 86,54\%. Partial fine-tuning nằm giữa hai chế độ này.

Head-only giữ nguyên cách backbone mô tả ảnh và chỉ học bộ phân loại trên các đặc trưng đó. Khi chuyển sang partial hoặc full fine-tuning, mô hình còn được thay đổi các đặc trưng để phù hợp với CIFAR-100. Kết quả cho thấy bước thích nghi này có ích ở cả CNN và Transformer. Đặc biệt, Swin-T tăng thêm 8,20 điểm phần trăm khi chuyển từ partial sang full: cập nhật các stage trước cũng đem lại lợi ích, dù stage cuối đã được học ở chế độ partial.

ResNet18 scratch đạt 64,53\%, cao hơn head-only nhưng thấp hơn partial và full. Như vậy, dùng trọng số pretrained không tự bảo đảm kết quả tốt hơn nếu backbone bị giữ nguyên. Mức lợi ích phụ thuộc vào cả biểu diễn ban đầu và cách cho phép mô hình thích nghi với dữ liệu mới.

\subsection{Swin-T scratch trong lịch 20 epoch}

Swin-T scratch chỉ đạt accuracy test trung bình 2,19\% trong lịch 20 epoch. Điểm thấp không chỉ xuất hiện trên test: ở cuối lịch, hai seed có accuracy train dưới 1\% và cross-entropy gần $\ln(100)\approx4{,}61$. Đây là mức loss của dự đoán gần đều trên 100 lớp. Mô hình chưa học tốt ngay cả dữ liệu dùng để cập nhật trọng số, nên vấn đề chính ở các lượt này là quá trình học chưa hiệu quả.

Confusion matrix giúp thấy rõ hành vi đó. Ở checkpoint của seed 42, mô hình dự đoán mọi ảnh test vào cùng một lớp. Vì test có 100 ảnh cho mỗi lớp, accuracy vẫn bằng 1\%, nhưng F1 của phần lớn lớp bằng 0. Điều này giải thích vì sao macro-F1 trung bình của cấu hình chỉ đạt 0,0067, thấp hơn accuracy 0,0219 khi cùng viết trên thang 0--1.

ResNet18 dùng convolution để kết hợp thông tin cục bộ ngay từ những tầng đầu. Swin-T học cách kết hợp các patch bằng attention, còn bản pretrained đã có các trọng số học từ ImageNet. Trong lịch 20 epoch đang xét, hai cách khởi tạo Swin-T vì vậy bắt đầu ở mức biểu diễn rất khác nhau. Đối chứng 40 epoch kiểm tra lợi ích của pretraining trong điều kiện bản scratch đã học được, thay vì chỉ dựa vào khoảng cách với một cấu hình gần mức đoán ngẫu nhiên.

\subsection{Độ chính xác và chi phí tính toán}

Ở chế độ full fine-tuning, Swin-T cao hơn ResNet18 trung bình $7{,}61 \pm 0{,}17$ điểm phần trăm khi so sánh các cặp cùng seed. Swin-T tốt hơn ở cả ba cặp. Đổi lại, thời gian train dài gấp 3,21 lần và bộ nhớ tensor cấp phát cực đại lớn gấp 4,70 lần. Vì vậy, Swin-T full là lựa chọn tốt hơn về accuracy trong thí nghiệm này, còn ResNet18 full tiết kiệm tài nguyên hơn.

Đóng băng các tầng giúp giảm chi phí backward và lưu gradient, nhưng mô hình vẫn phải forward qua backbone để tạo đặc trưng. Điều này giải thích vì sao số tham số trainable giảm mạnh mà thời gian có thể giảm ít. Với ResNet18, ba chế độ pretrained đều mất khoảng 64--67 giây train. Với Swin-T, partial chỉ mất khoảng một nửa thời gian của full, nhưng accuracy thấp hơn 8,20 điểm phần trăm.

Stage cuối chứa nhiều kênh đặc trưng, nên partial vẫn cập nhật một tỷ lệ đáng kể tham số: khoảng ba phần tư mô hình ResNet18 và hơn một nửa Swin-T. Bảng chi phí cho thấy rõ sự khác nhau giữa số tham số tổng và số tham số được học. Vì vậy, số tham số trainable chỉ phản ánh một phần chi phí; accuracy, thời gian và bộ nhớ bổ sung những mặt khác nhau của phép so sánh.

\subsection{Khoảng cách train/validation và lỗi theo lớp}

Cuối lịch 20 epoch, cả hai bản full đều đạt accuracy train gần 99\%, trong khi validation thấp hơn rõ rệt. ResNet18 có khoảng cách lớn hơn Swin-T. Đường cong cho thấy mô hình tiếp tục khớp dữ liệu train nhưng khả năng phân loại ảnh validation tăng ít hơn, phù hợp với overfitting. Chỉ số train được đo trên ảnh augmentation trong lúc cập nhật, còn validation được đo sau epoch ở chế độ eval; vì vậy, khoảng cách này dùng để đọc xu hướng học. Việc chọn checkpoint theo validation giúp lấy mô hình ở thời điểm phân loại tốt nhất trên tập này, thay vì mặc nhiên lấy epoch cuối.

Lỗi còn tập trung ở các lớp gần nhau về ngoại hình. Khi gộp ba seed, các lớp có recall thấp nhất của ResNet18 full gồm \code{boy}, \code{girl} và \code{otter}; với Swin-T full là \code{girl}, \code{boy} và \code{oak_tree}. Bảng sau trình bày các giá trị này. Mỗi lớp có 300 dự đoán, gồm 100 ảnh test được đánh giá bởi từng seed.

\begin{table}[htbp]\centering\small
\begin{tabular}{lrlr}\toprule
\multicolumn{2}{c}{ResNet18 full} & \multicolumn{2}{c}{Swin-T full}\\
Lớp & Recall (\%) & Lớp & Recall (\%)\\\midrule
\code{boy} & 50,0 & \code{girl} & 61,7\\
\code{girl} & 53,0 & \code{boy} & 63,3\\
\code{otter} & 57,0 & \code{oak_tree} & 67,3\\\bottomrule
\end{tabular}
\caption{Ba lớp có recall thấp nhất của mỗi bản full fine-tuning, tính từ confusion matrix test gộp seed 42, 43 và 44; support 300 dự đoán mỗi lớp.}
\end{table}

Ở cả hai mô hình, cặp nhầm phổ biến nhất là \code{oak_tree} thành \code{maple_tree}. Swin-T mắc 71 lỗi loại này, nhiều hơn 56 lỗi của ResNet18, dù accuracy toàn bộ cao hơn. Kết quả cho thấy một mô hình tốt hơn về điểm tổng thể vẫn có thể yếu hơn trên một cặp lớp cụ thể. Các ảnh sai có confidence cao trong gallery minh họa sự giống nhau về hình dáng và kết cấu ở độ phân giải $32\times32$; chúng được chọn để xem lỗi khó, không phải để ước lượng tỷ lệ từng loại lỗi.
